% GENERATED FILE. Edit canonical lessons, then run npm run generate:books.
% canonical-source-hash: fnv1a64:4f222d9009451a3e
% canonical-lessons: PT-C444-poste, PT-C444-ecoponto, PT-C444-sarjeta, PT-C444-pilarete, PT-C444-obstruir

\chapter{Coisas da rua}
\label{ch:pt-coisas-da-rua}
\hlchaptermodality{444}

\begin{chapteropening}
\textbf{By the end of this chapter:} \emph{I can describe things on the street: a post, the recycling bins, the gutter, the bollard and blocking the way.}

Complete the last lesson of chapter 444: I can describe things on the street: a post, the recycling bins, the gutter, the bollard and blocking the way.
\end{chapteropening}

\section[o poste]{o poste — the post, the pole}

\label{lesson:PT-C444-poste}



\begin{quote}
Before the new one: say the Portuguese for to report, then the Portuguese for to apply for.
\end{quote}

\subsection*{You'll want to know: o poste}
\textbf{o poste} — \textquotedblleft{}the post, the pole\textquotedblright{}.

\begin{quote}
O carro bateu num poste de luz.
\end{quote}

\begin{quote}
\emph{The car hit a lamp post.}
\end{quote}

\emph{O poste de luz} — the lamp post.

\begin{cousinweb}
\emph{poste} goes back to Latin \emph{postis} 'doorpost' (Lewis \& Short), probably through French \emph{poste}.
\end{cousinweb}

\begin{grammarlens}[title={What you've built}]
One more word for telling, explaining and describing at B1.
\end{grammarlens}

\subsection*{Your turn}
\begin{itemize}
\raggedright
  \item \emph{Say it:} \emph{o poste}
  \item \emph{Say it:} \emph{o poste}, once more
  \item \emph{Say it:} a sentence of your own with \emph{o poste}
  \item \emph{Recall:} say the Portuguese for to report, then the Portuguese for to apply for, then say \emph{o poste} again
\end{itemize}

\begin{tcolorbox}[breakable,colback=teal!4,colframe=teal!35!black,title={Before you move on}]
What does o poste mean? (\textquotedblleft{}The post, the pole\textquotedblright{}.) Say it once more.
\end{tcolorbox}
\section[o ecoponto]{o ecoponto — the recycling bins, the recycling point}

\label{lesson:PT-C444-ecoponto}



\begin{quote}
Before the new one: say the Portuguese for to apply for, then the Portuguese for the post.
\end{quote}

\subsection*{You'll want to know: o ecoponto}
\textbf{o ecoponto} — \textquotedblleft{}the recycling bins, the recycling point\textquotedblright{}.

\begin{quote}
Levamos as garrafas ao ecoponto.
\end{quote}

\begin{quote}
\emph{We take the bottles to the recycling bins.}
\end{quote}

Green for glass, blue for paper, yellow for plastic and metal.

\begin{cousinweb}
\emph{ecoponto} is a modern formation from \emph{eco-}, cut from \emph{ecologia}, and \emph{ponto} 'point', inherited from Latin \emph{pūnctum} (Lewis \& Short).
\end{cousinweb}

\begin{grammarlens}[title={What you've built}]
One more word for telling, explaining and describing at B1.
\end{grammarlens}

\subsection*{Your turn}
\begin{itemize}
\raggedright
  \item \emph{Say it:} \emph{o ecoponto}
  \item \emph{Say it:} \emph{o ecoponto}, once more
  \item \emph{Say it:} a sentence of your own with \emph{o ecoponto}
  \item \emph{Recall:} say the Portuguese for to apply for, then the Portuguese for the post, then say \emph{o ecoponto} again
\end{itemize}

\begin{tcolorbox}[breakable,colback=teal!4,colframe=teal!35!black,title={Before you move on}]
What does o ecoponto mean? (\textquotedblleft{}The recycling bins, the recycling point\textquotedblright{}.) Say it once more.
\end{tcolorbox}
\section[a sarjeta]{a sarjeta — the gutter (drain at the kerb)}

\label{lesson:PT-C444-sarjeta}



\begin{quote}
Before the new one: say the Portuguese for the post, then the Portuguese for the recycling bins.
\end{quote}

\subsection*{You'll want to know: a sarjeta}
\textbf{a sarjeta} — \textquotedblleft{}the gutter (drain at the kerb)\textquotedblright{}.

\begin{quote}
A chuva corre pela sarjeta.
\end{quote}

\begin{quote}
\emph{The rain runs along the gutter.}
\end{quote}

Figurative: \emph{cair na sarjeta} — to hit rock bottom.

\begin{cousinweb}
The origin of \emph{sarjeta} is not settled, and EtymDB has no entry for it.
\end{cousinweb}

\begin{grammarlens}[title={What you've built}]
One more word for telling, explaining and describing at B1.
\end{grammarlens}

\subsection*{Your turn}
\begin{itemize}
\raggedright
  \item \emph{Say it:} \emph{a sarjeta}
  \item \emph{Say it:} \emph{a sarjeta}, once more
  \item \emph{Say it:} a sentence of your own with \emph{a sarjeta}
  \item \emph{Recall:} say the Portuguese for the post, then the Portuguese for the recycling bins, then say \emph{a sarjeta} again
\end{itemize}

\begin{tcolorbox}[breakable,colback=teal!4,colframe=teal!35!black,title={Before you move on}]
What does a sarjeta mean? (\textquotedblleft{}The gutter (drain at the kerb)\textquotedblright{}.) Say it once more.
\end{tcolorbox}
\section[o pilarete]{o pilarete — the bollard}

\label{lesson:PT-C444-pilarete}



\begin{quote}
Before the new one: say the Portuguese for the recycling bins, then the Portuguese for the gutter.
\end{quote}

\subsection*{You'll want to know: o pilarete}
\textbf{o pilarete} — \textquotedblleft{}the bollard\textquotedblright{}.

\begin{quote}
Há pilaretes para os carros não subirem ao passeio.
\end{quote}

\begin{quote}
\emph{There are bollards to stop cars driving onto the pavement.}
\end{quote}

A small form of \emph{o pilar}, the pillar.

\begin{cousinweb}
\emph{pilarete} is a small form of \emph{pilar} 'pillar', which goes back to Latin \emph{pīla} 'pillar' (Lewis \& Short).
\end{cousinweb}

\begin{grammarlens}[title={What you've built}]
One more word for telling, explaining and describing at B1.
\end{grammarlens}

\subsection*{Your turn}
\begin{itemize}
\raggedright
  \item \emph{Say it:} \emph{o pilarete}
  \item \emph{Say it:} \emph{o pilarete}, once more
  \item \emph{Say it:} a sentence of your own with \emph{o pilarete}
  \item \emph{Recall:} say the Portuguese for the recycling bins, then the Portuguese for the gutter, then say \emph{o pilarete} again
\end{itemize}

\begin{tcolorbox}[breakable,colback=teal!4,colframe=teal!35!black,title={Before you move on}]
What does o pilarete mean? (\textquotedblleft{}The bollard\textquotedblright{}.) Say it once more.
\end{tcolorbox}
\section[obstruir]{obstruir — to obstruct, to block}

\label{lesson:PT-C444-obstruir}



\begin{quote}
Before the new one: say the Portuguese for the gutter, then the Portuguese for the bollard.
\end{quote}

\subsection*{You'll want to know: obstruir}
\textbf{obstruir} — \textquotedblleft{}to obstruct, to block\textquotedblright{}.

\begin{quote}
Um camião obstruiu a rua durante a manhã.
\end{quote}

\begin{quote}
\emph{A truck blocked the street all morning.}
\end{quote}

\emph{Obstruir a passagem} — to block the way.

\begin{cousinweb}
\emph{obstruir} is a learned borrowing from Latin \emph{obstruere} 'to build up against, to block' (EtymDB), on \emph{struere} 'to build'.
\end{cousinweb}

\begin{grammarlens}[title={What you've built}]
One more word for telling, explaining and describing at B1.
\end{grammarlens}

\subsection*{Your turn}
\begin{itemize}
\raggedright
  \item \emph{Say it:} \emph{obstruir}
  \item \emph{Say it:} \emph{obstruir}, once more
  \item \emph{Say it:} a sentence of your own with \emph{obstruir}
  \item \emph{Recall:} say the Portuguese for the gutter, then the Portuguese for the bollard, then say \emph{obstruir} again
\end{itemize}

\begin{tcolorbox}[breakable,colback=teal!4,colframe=teal!35!black,title={Before you move on}]
What does obstruir mean? (\textquotedblleft{}To obstruct, to block\textquotedblright{}.) Say it once more.
\end{tcolorbox}
